\section{Core Module Implementation}

\subsection{Configuration Management Module (config.py)}

The configuration management module is responsible for loading and managing the
system's configuration parameters, ensuring that all modules use a unified
configuration.

\subsubsection{Function Implementation}
\begin{itemize}
    \item Loads configuration from `configs/config.yaml` file
    \item Provides a global singleton configuration object `GLOBAL\_CONFIG` for direct
          import by other modules
    \item Supports customization of configuration file paths
\end{itemize}

\subsubsection{Key Code}

\begin{lstlisting}[language=Python]
def load_config(config_path: str | Path = DEFAULT_CONFIG_PATH) -> Dict[str, Any]:
    if not os.path.exists(config_path):
        raise FileNotFoundError(f"Config file {config_path} not found.")
    try:
        with open(config_path, "r", encoding="utf-8") as f:
            config = yaml.safe_load(f)
            return config
    except yaml.YAMLError as e:
        raise ValueError(f"YAML error in {config_path}: {e}")

GLOBAL_CONFIG = load_config()
\end{lstlisting}

\subsection{Data Ingestion Module (ingest)}

The data ingestion module is responsible for parsing, chunking, and vectorizing
financial PDF reports, providing basic data for subsequent retrieval and
generation.

\subsubsection{PDF Parsing (pdf\_parser.py)}

\paragraph{Function Implementation}
\begin{itemize}
    \item Uses PyMuPDFReader to parse PDF reports and extract text content
    \item Ensures that each document's metadata contains `source` and `page\_label`
    \item Normalizes metadata format to ensure compliance with the system's data contract
\end{itemize}

\paragraph{Key Code}
\begin{lstlisting}[language=Python]
def load_financial_pdfs(data_dir: str | Path = None) -> List[Document]:
    parser = PyMuPDFReader()
    for pdf_path in pdf_files:
        file_name = pdf_path.name
        try:
            docs_from_file = parser.load_data(file_path=str(pdf_path))
            for i, doc in enumerate(docs_from_file):
                raw_metadata = doc.metadata or {}
                extracted_page = raw_metadata.get("page", raw_metadata.get("source_page", str(i + 1)))
                doc.metadata = {
                    "source": file_name,
                    "page_label": str(extracted_page),
                    "total_pages": raw_metadata.get("total_pages", "unknown"),
                }
                doc.excluded_embed_metadata_keys = ["page_label", "total_pages", "source"]
                doc.excluded_llm_metadata_keys = ["total_pages"]
                all_documents.append(doc)
        except Exception as e:
            print(f"Failed: {file_name} - {str(e)}")
    return all_documents
\end{lstlisting}

\subsubsection{Chunking Strategy (chunker.py)}

\paragraph{Function Implementation}
\begin{itemize}
    \item Implements baseline chunking strategy: uses fixed-size chunks with a chunk size
          of 256 tokens
    \item Implements semantic chunking strategy: based on BGE-M3 semantic chunking, which
          splits by calculating cosine similarity between sentences
\end{itemize}

\subsubsection{Index Construction (indexer.py)}

\paragraph{Function Implementation}
\begin{itemize}
    \item Uses BGE-M3 model to vectorize chunked text
    \item Stores vectors in local ChromaDB
    \item Supports RAPTOR tree-based summary construction
\end{itemize}

\subsection{Retrieval Module (retrieval)}

The retrieval module is responsible for retrieving relevant information from
the vector database based on user queries, providing context for the generation
module.

\subsubsection{Retrieval Strategy (retriever.py)}

\paragraph{Function Implementation}
\begin{itemize}
    \item Implements single-path dense vector retrieval (Phase 1)
    \item Implements multi-path hybrid retrieval (Phase 2): integrates BM25 retrieval and
          vector retrieval, fuses results through RRF
    \item Implements reranking: uses bge-reranker-base to rerank retrieval results
    \item Implements global singleton cache to improve performance
\end{itemize}

\paragraph{Key Code}
\begin{lstlisting}[language=Python]
def get_hybrid_retriever(index: VectorStoreIndex) -> BaseRetriever:
    global _CACHED_BM25_RETRIEVER
    vector_top_k = GLOBAL_CONFIG["retrieval"].get("vector_top_k", 20)
    bm25_top_k = GLOBAL_CONFIG["retrieval"].get("bm25_top_k", 20)
    vector_retriever = index.as_retriever(similarity_top_k=vector_top_k)
    if _CACHED_BM25_RETRIEVER is None:
        print("Building BM25 retriever...")
        nodes_for_bm25 = _extract_nodes_from_vector_store(index)
        if not nodes_for_bm25:
            return vector_retriever
        _CACHED_BM25_RETRIEVER = BM25Retriever.from_defaults(
            nodes=nodes_for_bm25, similarity_top_k=bm25_top_k
        )
    fusion_retriever = QueryFusionRetriever(
        retrievers=[vector_retriever, _CACHED_BM25_RETRIEVER],
        similarity_top_k=GLOBAL_CONFIG["retrieval"].get("vector_top_k", 20),
        num_queries=1,
        mode="reciprocal_rerank",
        llm=Ollama(
            model=GLOBAL_CONFIG["llm"]["weak_model"],
            base_url=GLOBAL_CONFIG["llm"]["ollama_base_url"],
            temperature=0.0,
            request_timeout=300.0,
        ),
    )
    return fusion_retriever
\end{lstlisting}

\subsubsection{Reranking Module (reranker.py)}

\paragraph{Function Implementation}
\begin{itemize}
    \item Uses bge-reranker-base to rerank retrieval results
    \item Implements global singleton cache to improve performance
\end{itemize}

\subsection{Generation Module (generation)}

The generation module is responsible for building the dialogue engine and
generating answers, providing intelligent financial report analysis for users.

\subsubsection{Dialogue Engine (pipeline.py)}

\paragraph{Function Implementation}
\begin{itemize}
    \item Uses CondensePlusContextChatEngine as the core dialogue engine
    \item Integrates ChatMemoryBuffer to implement short-term memory management
    \item Supports multi-turn conversations and understands contextual references
\end{itemize}

\subsubsection{LLM Backend (llm\_backend.py)}

\paragraph{Function Implementation}
\begin{itemize}
    \item Encapsulates locally deployed large language models via Ollama
    \item Supports switching between different models
\end{itemize}

\subsubsection{Prompt Design (prompt.py)}

\paragraph{Function Implementation}
\begin{itemize}
    \item Designs strict anti-hallucination and mandatory citation instructions
    \item Ensures generated answers contain citation information
\end{itemize}

\subsection{Evaluation Module (evaluation)}

The evaluation module is responsible for evaluating the system's performance to
ensure the quality of system output.

\subsubsection{RAGAS Evaluation (ragas\_eval.py)}

\paragraph{Function Implementation}
\begin{itemize}
    \item Uses RAGAS evaluation framework to evaluate system answers
    \item Registers locally running Qwen2.5-72B as RAGAS's judge\_model
    \item Calculates Context Precision, Faithfulness, and Answer Relevancy metrics
\end{itemize}

\subsubsection{Citation Audit (citation\_audit.py)}

\paragraph{Function Implementation}
\begin{itemize}
    \item Audits citation information in generated answers
    \item Ensures cited page numbers and source files exist
\end{itemize}

\subsection{Interaction Module (ui)}

The interaction module is responsible for building the user interface and
providing a friendly user experience.

\subsubsection{Frontend Interface (app.py)}

\paragraph{Function Implementation}
\begin{itemize}
    \item Uses Gradio to build the frontend interface
    \item Implements multi-document management, conversational question-answering,
          citation jump, and artifact generation functions
    \item Supports clicking citations to jump to corresponding pages in PDF
\end{itemize}
